\documentclass[10pt,a4paper]{amsart}
\usepackage[margin=19mm]{geometry}
\usepackage{amsmath,amssymb,mathtools}
\usepackage[scr=rsfs,cal=euler]{mathalfa}
\usepackage{xfrac}
\usepackage[unicode,hidelinks,bookmarks=false]{hyperref}
\newcommand{\ie}{i.e.\ }
\newcommand{\cf}{cf.\ }
\newcommand{\resp}{respectively\ }
\newcommand{\Subsection}{\subsection}
\providecommand{\nobreakdash}{\nobreak\hbox{-}}
\setlength{\emergencystretch}{2em}
\multlinegap0pt
\usepackage{amssymb}
\usepackage{xcolor}
\newtheorem{lemma}{Lemma}
\newcommand{\Nb}{{\mathbb N}}
\newcommand{\cU}{{\mathcal U}}
\newcommand{\diam}{{\rm diam}}
\newcommand{\mesh}{{\rm mesh}}
\title{On the denseness of distal points}
\author{Su Gao, Hanfeng Li, Ruiwen Li}
\date{Source: arXiv:2607.26446v1 (CC BY 4.0)}
\begin{document}
\maketitle

\noindent\emph{Source and licence.} On the denseness of distal points. Version arXiv:2607.26446v1 (CC BY 4.0), distributed by arXiv under the Creative Commons Attribution 4.0 International licence, http://creativecommons.org/licenses/by/4.0/, as recorded in the arXiv licence metadata for this exact version and shown on the abstract page https://arxiv.org/abs/2607.26446v1. Full licence notice: \texttt{licenses/arxiv-2607.26446-CC-BY-4.0.txt}.\par\smallskip

\noindent\textit{Editorial scope.} Counted unit: the article's lemma orbit (source0.tex lines 718--726). The article's complete original proof of it is delivered with it (source0.tex lines 727--758, one \texttt{proof} environment of 32 lines). No mathematics is written, repaired or re-derived: the statement and the proof are the article's own lines, in the article's own notation, and no retained line is substituted or re-flowed.  No further source-local context is needed: every cross-reference of every retained line resolves inside the two delivered spans.  Nothing was omitted from any retained span, so the package carries no omission record.  No retained line contains a citation, so the wrapper carries no bibliography and no bibliography entry.  No retained line contains a figure environment, an image-inclusion command or drawing code, so no image is delivered and the package declares no asset dependency.  Editorial formatting only, in addition: an AMS \texttt{amsart} preamble that restates the article's own declarations and macros used by the retained lines, namely the environments \texttt{lemma} on separate counters, together with the standard packages those lines need.  The retained environments print in this document as Lemma 1 in order of appearance, whereas the article prints the same items with the numbering of its own full document.  The complete native source of the article as distributed in the arXiv e-print for this exact version is bundled unchanged as \texttt{sources/206269/source0.tex}.
\begin{lemma} \label{L-avoid orbit}
Let $G$ act on a compact metrizable space $X$ continuously. Let $x\in X$ with dense orbit. Let $\rho$ be a compatible metric on $X$. Then there is a sequence $\{\cU_n\}_{n\in \Nb}$ of finite families of nonempty open subsets of $X$ such that the following hold:
\begin{enumerate}
\item[(a)] for each $n\in \Nb$, the members of $\cU_n$ are pairwise disjoint and the union $\bigcup \cU_n$ is dense in $X$;
\item[(b)] for each $n\in \Nb$, every member of $\cU_{n+1}$ is contained in some member of $\cU_n$;
\item[(c)] ${\rm mesh}(\cU_n):=\max_{U\in \cU_n}\diam(U)$ converges to $0$ as $n\to \infty$;
\item[(d)] for each $n\in \Nb$, the orbit of $x$ is contained in $\bigcup \cU_n$.
\end{enumerate}
\end{lemma}

\begin{proof} The lemma is trivial when $X$ is finite. Thus we may assume that $X$ is infinite. We shall construct $\{\cU_n\}_{n\in \Nb}$ inductively on $n$ satisfying conditions (a), (b), (d) and the following condition which is stronger than condition (c):
\begin{enumerate}
\item[(c')] for each $n\in \Nb$,  ${\rm mesh}(\cU_n)\le \diam(X)/2^{n-1}$. 
\end{enumerate}

To start, we take $\cU_1=\{X\}$. Assume that we have constructed $\cU_1, \dots, \cU_n$ for some $n\in \Nb$. Since $X$ is compact, for each $U\in \cU_n$ we can find a finite subset $A_U$ of $U$ such that for every $x\in U$, there is some $a\in A_U$ with $\rho(x, a)<r_n:=\diam(X)/2^{n+1}$. Set $A:=\bigcup_{U\in \cU_n}A_U$. Since $X$ is infinite and the orbit of $x$ is dense in $X$, the orbit of $x$ is infinite. Enumerate the points in the orbit of $x$ as $x_1, x_2, \dots$. 

For any $y\in X$ and $r>0$, we write $B(y, r)$ for the open ball $\{z\in X: \rho(z, y)<r\}$. 
We shall construct open sets $V_{a, j}$ for  $a\in A$ and $j\in \Nb$ such that the following hold:
\begin{enumerate}
\item[(i)] for each $U\in \cU_n$ and $a\in A_U$, one has $V_{a, j}\subseteq V_{a, j+1}$ and $\overline{V_{a, j}}\subseteq U\cap B(a, r_n)$ for all $j\in \Nb$;
\item[(ii)] for each $j\in \Nb$, the sets $\overline{V_{a, j}}$ for $a\in A$ are pairwise disjoint, and $x_j\in \bigcup_{a\in A} V_{a, j}$. 
\end{enumerate}
Assume that we have constructed such $V_{a, j}$. Set $V_a=\bigcup_{j\in \Nb}V_{a, j}$ for $a\in A$, and $\cU_{n+1}=\{V_a: a\in A, V_a\neq \varnothing\}$. Then the members of $\cU_{n+1}$ are pairwise disjoint. Since the orbit of $x$ is dense in $X$, and the orbit of $x$ is contained in $\bigcup \cU_{n+1}$, we see that $\bigcup \cU_{n+1}$ is dense in $X$. For each $U\in \cU_n$ and $a\in A_U$, one has $V_a\subseteq U\cap B(a, r_n)$. Note that $\mesh(\cU_{n+1})\le 2\cdot r_n=\diam(X)/2^n$. Thus $\cU_1, \dots, \cU_{n+1}$ satisfy conditions (a), (b), (c') and (d).

Now it suffices to construct open sets $V_{a, j}$ satisfying (i) and (ii). We construct them inductively on $j$. 

Let $j=1$. By condition (d), the point $x_1$ is in some $U\in \cU_n$. Then we can find some $a\in A_U$ such that $x_1\in B(a, r_n)$. Take a small $r>0$ with $\overline{B(x_1, r)}\subseteq U\cap B(a, r_n)$. Setting $V_{a, 1}=B(x_1, r)$ and $V_{b, 1}=\varnothing$ for all $b\in A\setminus \{a\}$ satisfies the requirement. 

Assume that we have constructed $V_{a, j}$ for some $j\in \Nb$ and all $a\in A$. 
We separate the discussion into 3 cases. 

Case I: $x_{j+1}\in \bigcup_{b\in A} V_{b, j}$. In this case we may set $V_{b, j+1}=V_{b, j}$ for all $b\in A$. 

Case II: $x_{j+1}\not\in \bigcup_{b\in A} \overline{V_{b, j}}$. By condition (d), the point $x_{j+1}$ is in some $U\in \cU_n$. Then we can find some $a\in A_U$ such that $x_{j+1}\in B(a, r_n)$. Take a small $r>0$ such that  $\overline{B(x_{j+1}, r)}\subseteq U\cap B(a, r_n)$ and 
$\overline{B(x_{j+1}, r)}$ is disjoint with $\bigcup_{b\in A} \overline{V_{b, j}}$. Setting $V_{a, j+1}=B(x_{j+1}, r)\cup V_{a, j}$ and $V_{b, j+1}=V_{b, j}$ for all $b\in A\setminus \{a\}$ satisfies the requirement. 

Case III: $x_{j+1}\in \partial V_{a, j}$ for some $a\in A$. Then $a$ is in some $U\in \cU_n$. By the condition (ii) we have $x_{j+1}\not\in \overline{V_{b, j}}$ for all $b\in A\setminus \{a\}$ and $x_{j+1}\in U\cap B(a, r_n)$. Take a small $r>0$ such that  $\overline{B(x_{j+1}, r)}\subseteq U\cap B(a, r_n)$ and 
$\overline{B(x_{j+1}, r)}$ is disjoint with $\bigcup_{b\in A\setminus \{a\}} \overline{V_{b, j}}$. Setting $V_{a, j+1}=B(x_{j+1}, r)\cup V_{a, j}$ and $V_{b, j+1}=V_{b, j}$ for all $b\in A\setminus \{a\}$ satisfies the requirement. 

This finishes the construction of $V_{a, j+1}$ for $a\in A$.
\end{proof}

\end{document}
